\section{Résultats phytosanitaires}
\subsection{Maladies fongiques observées et efficacité des traitements appliqués}
Les maladies fongiques observées durant l'expérimentation sont la \textbf{fusariose, le phytophthora et l’anthracnose}. Comme souligné au chapitre 3, nous avons utilisé des traitements comme la solution d’ail et la solution de bicarbonate dont les effets diffèrent.

Après utilisation de la \textbf{solution de bicarbonate}, nous avons observé des \textbf{brûlures sur les feuilles de gombo}, qui s’aggravaient sous l’effet du soleil dans la journée. Ainsi, nous avons procédé par rinçage afin de stopper ou de réduire les dégâts du bicarbonate. Par contre, \textbf{la solution d’ail} avait un \textbf{effet positif mais pas très efficace pour soigner} les plantes déjà affectées.

\begin{figure}[H]
    \centering
    \imgstd{images/z14_gombo_bruler_chap5.jpg}
    \caption{Action de la solution de bicarbonate sur les feuilles de gombo. }
    \label{fig:gombo-bruler}
\end{figure}
